% GENERATED FILE. Edit canonical lessons, then run npm run generate:books.
% canonical-source-hash: fnv1a64:d343051c513e03c8
% canonical-lessons: HI-C162-tairna, HI-C162-nacna, HI-C162-hasna, HI-C162-rona, HI-C162-bhulna

\chapter{To swim, To dance, To laugh, To cry, To forget}
\label{ch:hi-to-swim-to-dance-to-laugh-to-cry-to-forget}
\hlchaptermodality{162}

\begin{chapteropening}
\textbf{By the end of this chapter:} \emph{I can say to swim, to dance, to laugh, to cry and to forget.}

Complete the last lesson of chapter 162: I can say to swim, to dance, to laugh, to cry and to forget.
\end{chapteropening}

\section[\texorpdfstring{\hi{तैरना} (tairnā)}{तैरना (tairnā)}]{\hi{तैरना} (tairnā) — to swim}

\label{lesson:HI-C162-tairna}



\begin{quote}
Before the new one: say the Hindi for to fall, then the Hindi for to fly.
\end{quote}

\subsection*{You'll want to know: \hi{तैरना}}
\textbf{\hi{तैरना}} — \emph{tairnā} — \textquotedblleft{}to swim\textquotedblright{}.

\begin{cousinweb}
\textbf{\hi{तैरना}} (\emph{tairnā}) is a tadbhava, a word inherited from Sanskrit and reshaped by sound change. It comes from \emph{tirate}, \textquotedblleft{}crosses over\textquotedblright{} (CDIAL 5821). CDIAL rejects a Dravidian source proposed by Bloch, and explains the \emph{ai} as a blend with another verb.
\end{cousinweb}

\begin{grammarlens}[title={What you've built}]
One more word for this chapter.
\end{grammarlens}

\subsection*{Your turn}
\begin{itemize}
\raggedright
  \item \emph{Say it:} \emph{tairnā}
  \item \emph{Say it:} \emph{tairnā}, once more
  \item \emph{Say it:} say \emph{tairnā}
  \item \emph{Recall:} say the Hindi for to fall, then the Hindi for to fly, then say \emph{tairnā} again
\end{itemize}

\begin{tcolorbox}[breakable,colback=teal!4,colframe=teal!35!black,title={Before you move on}]
What does \hi{तैरना} mean? (\textquotedblleft{}To swim\textquotedblright{}.) Say it once more.
\end{tcolorbox}
\section[\texorpdfstring{\hi{नाचना} (nācnā)}{नाचना (nācnā)}]{\hi{नाचना} (nācnā) — to dance}

\label{lesson:HI-C162-nacna}



\begin{quote}
Before the new one: say the Hindi for to fly, then the Hindi for to swim.
\end{quote}

\subsection*{You'll want to know: \hi{नाचना}}
\textbf{\hi{नाचना}} — \emph{nācnā} — \textquotedblleft{}to dance\textquotedblright{}.

\begin{cousinweb}
\textbf{\hi{नाचना}} (\emph{nācnā}) is a tadbhava, a word inherited from Sanskrit and reshaped by sound change. It comes from \emph{nṛtyati}, \textquotedblleft{}dances\textquotedblright{} (CDIAL 7583).
\end{cousinweb}

\begin{grammarlens}[title={What you've built}]
One more word for this chapter.
\end{grammarlens}

\subsection*{Your turn}
\begin{itemize}
\raggedright
  \item \emph{Say it:} \emph{nācnā}
  \item \emph{Say it:} \emph{nācnā}, once more
  \item \emph{Say it:} say \emph{nācnā}
  \item \emph{Recall:} say the Hindi for to fly, then the Hindi for to swim, then say \emph{nācnā} again
\end{itemize}

\begin{tcolorbox}[breakable,colback=teal!4,colframe=teal!35!black,title={Before you move on}]
What does \hi{नाचना} mean? (\textquotedblleft{}To dance\textquotedblright{}.) Say it once more.
\end{tcolorbox}
\section[\texorpdfstring{\hi{हँसना} (hãsnā)}{हँसना (hãsnā)}]{\hi{हँसना} (hãsnā) — to laugh}

\label{lesson:HI-C162-hasna}



\begin{quote}
Before the new one: say the Hindi for to swim, then the Hindi for to dance.
\end{quote}

\subsection*{You'll want to know: \hi{हँसना}}
\textbf{\hi{हँसना}} — \emph{hãsnā} — \textquotedblleft{}to laugh\textquotedblright{}.

\begin{cousinweb}
\textbf{\hi{हँसना}} (\emph{hãsnā}) is a tadbhava, a word inherited from Sanskrit and reshaped by sound change. It comes from \emph{hasati}, \textquotedblleft{}laughs\textquotedblright{} (CDIAL 14021).
\end{cousinweb}

\begin{grammarlens}[title={What you've built}]
One more word for this chapter.
\end{grammarlens}

\subsection*{Your turn}
\begin{itemize}
\raggedright
  \item \emph{Say it:} \emph{hãsnā}
  \item \emph{Say it:} \emph{hãsnā}, once more
  \item \emph{Say it:} say \emph{hãsnā}
  \item \emph{Recall:} say the Hindi for to swim, then the Hindi for to dance, then say \emph{hãsnā} again
\end{itemize}

\begin{tcolorbox}[breakable,colback=teal!4,colframe=teal!35!black,title={Before you move on}]
What does \hi{हँसना} mean? (\textquotedblleft{}To laugh\textquotedblright{}.) Say it once more.
\end{tcolorbox}
\section[\texorpdfstring{\hi{रोना} (ronā)}{रोना (ronā)}]{\hi{रोना} (ronā) — to cry}

\label{lesson:HI-C162-rona}



\begin{quote}
Before the new one: say the Hindi for to dance, then the Hindi for to laugh.
\end{quote}

\subsection*{You'll want to know: \hi{रोना}}
\textbf{\hi{रोना}} — \emph{ronā} — \textquotedblleft{}to cry\textquotedblright{}.

\begin{cousinweb}
\textbf{\hi{रोना}} (\emph{ronā}) is a tadbhava, a word inherited from Sanskrit and reshaped by sound change. It comes from \emph{roditi}, \textquotedblleft{}weeps\textquotedblright{} (CDIAL 10840).
\end{cousinweb}

\begin{grammarlens}[title={What you've built}]
One more word for this chapter.
\end{grammarlens}

\subsection*{Your turn}
\begin{itemize}
\raggedright
  \item \emph{Say it:} \emph{ronā}
  \item \emph{Say it:} \emph{ronā}, once more
  \item \emph{Say it:} say \emph{ronā}
  \item \emph{Recall:} say the Hindi for to dance, then the Hindi for to laugh, then say \emph{ronā} again
\end{itemize}

\begin{tcolorbox}[breakable,colback=teal!4,colframe=teal!35!black,title={Before you move on}]
What does \hi{रोना} mean? (\textquotedblleft{}To cry\textquotedblright{}.) Say it once more.
\end{tcolorbox}
\section[\texorpdfstring{\hi{भूलना} (bhūlnā)}{भूलना (bhūlnā)}]{\hi{भूलना} (bhūlnā) — to forget}

\label{lesson:HI-C162-bhulna}



\begin{quote}
Before the new one: say the Hindi for to laugh, then the Hindi for to cry.
\end{quote}

\subsection*{You'll want to know: \hi{भूलना}}
\textbf{\hi{भूलना}} — \emph{bhūlnā} — \textquotedblleft{}to forget\textquotedblright{}.

\begin{cousinweb}
\textbf{\hi{भूलना}} (\emph{bhūlnā}) is a tadbhava, a word inherited and reshaped by sound change. It comes from \emph{*bhull}, \textquotedblleft{}err, forget\textquotedblright{}, a form CDIAL reconstructs from its descendants (CDIAL 9538).
\end{cousinweb}

\begin{grammarlens}[title={What you've built}]
One more word for this chapter.
\end{grammarlens}

\subsection*{Your turn}
\begin{itemize}
\raggedright
  \item \emph{Say it:} \emph{bhūlnā}
  \item \emph{Say it:} \emph{bhūlnā}, once more
  \item \emph{Say it:} say \emph{bhūlnā}
  \item \emph{Recall:} say the Hindi for to laugh, then the Hindi for to cry, then say \emph{bhūlnā} again
\end{itemize}

\begin{tcolorbox}[breakable,colback=teal!4,colframe=teal!35!black,title={Before you move on}]
What does \hi{भूलना} mean? (\textquotedblleft{}To forget\textquotedblright{}.) Say it once more.
\end{tcolorbox}
